%& -shell-escape
% !TeX root = thesis-abstract.tex
% !TeX encoding = UTF-8
% !TeX program = XeLaTeX

\documentclass{.local/lib/classes/ndvr-super-article-standalone}

\begin{document}
    
    Временные моменты аномалий в двух дубликатах 
    видео могут не совпадать.
    Причинами этого могут быть особенности 
    временных систем координат в каждом видео и ошибки 
    при определении моментов аномалий.  

    Для решения этой проблемы предлагается 
    использовать относительные длины съёмок.
    Относительная длина съёмки вычисляется 
    как вектор отношений абсолютной длины
    съёмки к~абсолютным длинам остальных съёмок видео.
    В~практических задачах удобнее вычислять отношения длин
    для~трех предыдущих съёмок, а~не~для~всех.
    Это удобно и~в случае, если всё видео целиком недоступно,
    как например в~задачах реального времени.
    Относительные длины съёмок двух нечетких дубликатов 
    тоже могут не совпадать.
        
\end{document}
